\phantomsection

本人声明，在撰写题目为\textbf{《\varTitle》}的本科毕业论文过程中，对 AI 工具的使用情况如下：

\section*{一、使用的 AI 工具及使用目的}

\noindent
\begin{tabularx}{\textwidth}{|c|X|X|}
    \hline
    \textbf{序号} & \textbf{AI 工具} & \textbf{使用目的} \\
    \hline
    1 & Cursor AI 编程助手 & 辅助论文文本润色与病句检查、LaTeX 排版问题排查与格式调整、公开资料检索思路整理、图表说明文字修改 \\
    \hline
\end{tabularx}

\section*{二、使用程度说明}

\noindent
\begin{tabularx}{\textwidth}{|c|X|X|X|X|}
    \hline
    \textbf{序号} & \textbf{章节} & \textbf{AI 工具} & \textbf{使用方式} & \textbf{处理方式} \\
    \hline
    1 & 第1章 绪论、第2章 技术综述 & Cursor AI 编程助手 & 输入特定段落请其提出语言优化建议，获取多种表述方案 & 本人对 AI 提供的内容进行筛选、修改、重新组织语言，结合研究内容进行调整，确保学术表述准确 \\
    \hline
    2 & 第3章 系统设计、第4章 实现方案 & Cursor AI 编程助手 & 辅助梳理系统架构描述逻辑、代码注释与流程图说明文字 & 本人对建议内容进行核实后采纳，确保技术描述与实际实现一致 \\
    \hline
    3 & 第5章 实验评估 & Cursor AI 编程助手 & 辅助润色实验结果描述语句 & 本人结合实际实验数据对表述进行核实与修改，确保结论描述与数据一致 \\
    \hline
\end{tabularx}

\vspace{1em}
在论文的第1章绪论、第2章技术综述中使用了上述 AI 工具。使用方式为输入特定段落请其提出语言优化建议，获取多种表述方案。本人对 AI 工具提供的内容进行了筛选、修改、重新组织语言，结合自身研究内容进行调整，以确保论文内容符合学术要求及本人的研究思路。

\section*{三、原创性与责任声明}

本人确认，尽管使用了上述 AI 工具，但论文整体的研究思路、核心观点、论证过程及最终结论均为本人独立思考与研究的成果。论文中引用 AI 工具生成内容之处，均已按照学校规定的学术规范进行了明确标注。

若论文因 AI 工具使用不当或未按要求标注等原因出现学术不端问题，本人愿意承担全部责任。
